\chapter{Mô hình biểu diễn tri thức}\label{chap:model}
\section{Vận dụng mô hình COKB}
Mô hình COKB tổ chức tri thức theo sáu thành phần
\begin{equation}
  \K=(C,H,R,\mathrm{Ops},\mathrm{Funcs},\mathrm{Rules}).
  \label{eq:cokb}
\end{equation}
Trong lý thuyết, $C$ chứa các lớp đối tượng tính toán; $H$ mô tả phân cấp;
$R$ là các quan hệ khác; $\mathrm{Ops}$ và $\mathrm{Funcs}$ mô tả phép toán
và hàm; $\mathrm{Rules}$ chứa các luật suy diễn~\cite{do2015}.

Đồ án vận dụng cách tổ chức này cho bài toán ánh xạ mã.
Các đối tượng Python hiện chủ yếu là cấu trúc dữ liệu bất biến; các hành vi
đánh giá được đặt trong operators, rules và reasoner.
Vì vậy, đây là \emph{một hiện thực chuyên biệt theo cấu trúc COKB};
chưa hiện thực đầy đủ cơ chế đối tượng tự giải, giải phương trình nội bộ,
hay suy diễn đến bao đóng của một hệ COKB tổng quát.

\begin{table}[htbp]
\centering\small
\caption{Đối chiếu sáu thành phần COKB với source code}\label{tab:cokb}
\begin{tabularx}{\textwidth}{lYY}
\toprule
Thành phần & Cụ thể hóa trong đồ án & Hiện thực và giới hạn\\
\midrule
$C$ & CodeConcept, MappingAssertion, MappingAssessment, ProofTrace. &
\src{model.py} và vocabulary Turtle; một số lớp chỉ có trong schema.\\
$H$ & Phân cấp hệ mã, loại assessment, loại target. &
\src{rdfs:subClassOf}; không phải cây phân cấp bệnh dùng để suy luận nhãn.\\
$R$ & mapsFrom, mapsTo, assessesMapping, hasPremise. &
Quan hệ RDF do \src{repository.py} xuất.\\
Ops & Chuẩn hóa mã/nhãn, tìm xung đột, so sánh tập từ. &
\src{operators.py}; khử trùng cặp trong \src{importer.py}.\\
Funcs & Tìm mapping, đánh giá, giải thích, xuất kho tri thức. &
\src{find_mappings}, \src{infer}, \src{explain}, \src{save}.\\
Rules & Năm luật có ưu tiên và vết suy luận. &
\src{rules.py}, \src{reasoner.py}; logic chạy bằng Python.\\
\bottomrule
\end{tabularx}
\end{table}

Tên file Python trong bảng nằm dưới \src{scripts/cokb/}.
Ontology nằm ở \src{ontology/biomedical_mapping_cokb.ttl}.
Các cá thể Operator, Function và InferenceRule trong Turtle đóng vai trò
danh mục tri thức; chúng không tự thực thi được. Danh mục
\src{rules/mapping_rules.json} cũng không phải ngôn ngữ luật mà reasoner
đọc để chạy. Thứ tự thực thi thật được quyết định bởi \src{DEFAULT_RULES}.

\section{Các đối tượng và quan hệ}
\src{CodeConcept} gồm hệ mã, mã và nhãn. \src{MappingAssertion} liên kết hai
CodeConcept với bằng chứng nguồn. \src{MappingAssessment} bổ sung mức đánh
giá, loại tác nhân, lời giải thích và proof. \src{Premise} chứa bộ ba
thuộc tính \texttt{predicate}, \texttt{value}, \texttt{source}.
Những thuộc tính này thể hiện cụ thể tri thức mà phép đánh giá được phép dùng.

\begin{figure}[htbp]
\centering
\begin{tikzpicture}[node distance=0.9cm and 1.4cm]
  \node[block] (mapping) {MappingAssertion};
  \node[block,above left=of mapping,text width=2.8cm] (source) {CodeConcept\\nguồn};
  \node[block,above right=of mapping,text width=2.8cm] (target) {CodeConcept\\đích};
  \node[block,below=of mapping] (assessment) {MappingAssessment\\A / B / C / UNASSESSED};
  \node[block,below=of assessment] (proof) {ProofTrace\\luật và bằng chứng};
  \node[block,right=of proof,text width=2.5cm] (premise) {Premise};
  \draw[flow] (mapping) -- node[left,font=\scriptsize]{mapsFrom} (source);
  \draw[flow] (mapping) -- node[right,font=\scriptsize]{mapsTo} (target);
  \draw[flow] (assessment) -- node[right,font=\scriptsize]{assessesMapping} (mapping);
  \draw[flow] (proof) -- node[right,font=\scriptsize]{provesAssessment} (assessment);
  \draw[flow] (proof) -- node[above,font=\scriptsize]{hasPremise} (premise);
\end{tikzpicture}
\caption{Các đối tượng trung tâm và chiều quan hệ được xuất trong RDF.}
\label{fig:objects}
\end{figure}

Điểm đáng chú ý ở Hình~\ref{fig:objects} là mức A/B/C thuộc về
\emph{assessment của mapping}. Dữ kiện “mapping tồn tại” và kết luận
“mapping được đánh giá mức A” có chủ thể và trách nhiệm khác nhau.
Trong JSON, proof ghi \texttt{mapping\_id} cùng predicate và object;
trong RDF, liên kết \texttt{hasMappingLevel} xuất phát từ assessment.

\section{Phân cấp và mức độ hoàn thiện}
Ontology khai báo \src{BiomedicalCodeSystem} là lớp con của
\src{CodeSystem}; các lớp ICD nằm dưới \src{BiomedicalCodeSystem}.
Các lớp \src{RuleBasedAssessment}, \src{ModelAssessment},
\src{HumanAssessment} nằm dưới \src{MappingAssessment}.
Tuy nhiên, dữ liệu CodeConcept hiện ghi tên hệ mã bằng chuỗi
\src{codeSystem}; chưa liên kết tới một cá thể CodeSystem được quản lý phiên bản.

Tương tự, \src{SimpleTarget}, \src{OneOfTarget}, \src{AllOfTarget} đã có
trong vocabulary, nhưng \src{MappingAssertion} Python chỉ chứa một
\src{target}. Không thể suy từ việc có tên lớp trong ontology rằng
chương trình đã giải được lựa chọn “một trong nhiều mã” hoặc tổ hợp
“phải dùng đồng thời nhiều mã”.

\section{Dữ kiện nguồn, kết luận và bằng chứng}
Kho tri thức được xuất thành ba file:
\begin{equation}
  G_{\mathrm{export}}=G_{\mathrm{asserted}}\cup
  G_{\mathrm{inferred}}\cup G_{\mathrm{proof}}.
\end{equation}
\src{asserted_graph.ttl} lưu cặp mã, nhãn, tác nhân và nguồn.
\src{inferred_graph.ttl} lưu assessment, mức đánh giá và luật.
\src{proof_graph.ttl} lưu các premise cùng liên kết tới assessment.
Đây là ba tài liệu Turtle; lệnh xuất hiện tại không tự tạo ba named graph
trong một RDF Dataset.

Quan hệ \src{prov:wasDerivedFrom} được dùng để liên kết mapping với
EvidenceSource, phù hợp với mục đích mô tả nguồn gốc dữ liệu của
PROV-O~\cite{provo}. Ở bản cài đặt hiện tại, evidence là định danh chuỗi
do fixture hoặc tham số import cung cấp. Chưa có cơ chế xác thực nội dung
nguồn, snapshot nguồn, timestamp hay hash phiên bản luật trong proof.
Do đó, truy vết cấu trúc đã có nhưng khả năng kiểm toán lịch sử còn hạn chế.

\section{Ràng buộc tính hợp lệ}
Kiểm tra Python yêu cầu đủ hệ mã, mã, nhãn nguồn/đích và evidence source;
phát hiện trùng \texttt{mapping\_id}; cảnh báo ánh xạ một mã tới chính nó.
Kiểm tra SHACL áp dụng ràng buộc số lượng và kiểu liên kết trên dữ liệu
RDF~\cite{shacl}. Ví dụ, một assertion phải có đúng một \src{mapsFrom},
đúng một \src{mapsTo} và ít nhất một \src{prov:wasDerivedFrom}.

Gọi $V(M)$ là tập lỗi của một tập mapping $M$, trạng thái hợp lệ trong code là
\begin{equation}
  \operatorname{conforms}(M)
  \iff \nexists v\in V(M):\operatorname{severity}(v)=\texttt{error}.
\end{equation}
Một warning không làm trạng thái này chuyển thành false.
Tính hợp lệ về cấu trúc không chứng minh tính đúng của nhãn hay mapping.
Chẳng hạn, một chuỗi evidence sai nội dung nhưng không rỗng vẫn có thể
vượt qua kiểm tra hiện tại.
